% ═══════════════════════════════════════════════════════════════
% SLIDES: Oxide – Schrittweise Erarbeitung mit Funfacts
% Chemie Klasse 8 | Begleitmaterial zur Übung
% ═══════════════════════════════════════════════════════════════

\documentclass[aspectratio=169, 12pt]{beamer}

\usepackage[utf8]{inputenc}
\usepackage[T1]{fontenc}
\usepackage[ngerman]{babel}

\usepackage{helvet}
\renewcommand{\familydefault}{\sfdefault}

\usepackage{tikz}

% ═══════════════════════════════════════════════════════════════
% THEME
% ═══════════════════════════════════════════════════════════════

\usetheme{default}
\usecolortheme{default}

\definecolor{chemie}{HTML}{2a7a4b}
\definecolor{hellgrau}{HTML}{F0F0F0}
\definecolor{akzent}{HTML}{e65100}
\definecolor{loesung}{HTML}{C62828}

\setbeamercolor{frametitle}{fg=white, bg=chemie}
\setbeamercolor{structure}{fg=chemie}
\setbeamercolor{block title}{fg=white, bg=chemie}
\setbeamercolor{block body}{fg=black, bg=hellgrau}

\setbeamerfont{frametitle}{series=\bfseries}

\setbeamertemplate{navigation symbols}{}
\setbeamertemplate{footline}[frame number]

% Kompaktere Blöcke
\addtobeamertemplate{block begin}{}{\vspace{-2pt}}
\addtobeamertemplate{block end}{\vspace{-2pt}}{}

% ═══════════════════════════════════════════════════════════════
% BEFEHLE
% ═══════════════════════════════════════════════════════════════

\newcommand{\lsg}[1]{\textcolor{loesung}{\textbf{#1}}}
\newcommand{\luecke}{\underline{\hspace{2.5cm}}}

% ═══════════════════════════════════════════════════════════════
% DOKUMENT
% ═══════════════════════════════════════════════════════════════

\begin{document}

% ═══════════════════════════════════════════════════════════════
% TITEL
% ═══════════════════════════════════════════════════════════════

\begin{frame}
\begin{center}
\vspace{1.5cm}
{\Huge\bfseries\textcolor{chemie}{Oxide}}\\[1cm]
{\Large Formeln und Reaktionsgleichungen}\\[1.5cm]
{Chemie Klasse 8}
\end{center}
\end{frame}

% ═══════════════════════════════════════════════════════════════
% KREUZREGEL
% ═══════════════════════════════════════════════════════════════

\begin{frame}
\frametitle{Erinnerung: Kreuzregel}
\begin{center}
{\large Die Wertigkeit des einen wird zum Index des anderen.}
\vspace{0.8cm}

\begin{tikzpicture}[scale=1.3]
\node at (0,0) {\Huge Al};
\node at (0,-0.7) {(III)};
\node at (2.5,0) {\Huge O};
\node at (2.5,-0.7) {(II)};
\draw[->, thick, chemie] (0.3,-0.4) -- (2.2,0.2);
\draw[->, thick, chemie] (2.2,-0.4) -- (0.3,0.2);
\node at (5,0) {\Huge $\rightarrow$ \textcolor{chemie}{\textbf{Al$_2$O$_3$}}};
\end{tikzpicture}

\vspace{0.8cm}
{\large Gleiche Wertigkeiten $\rightarrow$ \textbf{kürzen!}}
\end{center}
\end{frame}

% ═══════════════════════════════════════════════════════════════
% 1. BARIUMOXID
% ═══════════════════════════════════════════════════════════════

\begin{frame}
\frametitle{1. Bariumoxid – Aufgabe}
\begin{block}{Gegeben}
Barium (Ba), Wertigkeit \textbf{II} \quad|\quad Sauerstoff (O), Wertigkeit \textbf{II}
\end{block}
\vspace{0.5cm}
\begin{block}{Aufgaben}
\begin{enumerate}
\item Oxidformel aufstellen: \luecke
\item Reaktionsgleichung (ausgeglichen): \quad \luecke{} + O$_2$ $\rightarrow$ \luecke
\end{enumerate}
\end{block}
\end{frame}

\begin{frame}
\frametitle{1. Bariumoxid – Lösung Schritt 1}
\begin{block}{Oxidformel}
Ba (II) + O (II) $\rightarrow$ Kreuzregel: 2:2 $\rightarrow$ \textbf{kürzen!} \hfill \lsg{BaO}
\end{block}
\vspace{0.3cm}
\begin{block}{Reaktionsgleichung (unausgeglichen)}
\lsg{Ba + O$_2$ $\rightarrow$ BaO}\\[0.3cm]
Links: 1 Ba, 2 O \quad Rechts: 1 Ba, 1 O \quad $\rightarrow$ \textbf{nicht ausgeglichen!}
\end{block}
\end{frame}

\begin{frame}
\frametitle{1. Bariumoxid – Lösung komplett}
\begin{center}
\vspace{0.5cm}
{\large Oxidformel:} \quad {\Huge \lsg{BaO}}
\vspace{1cm}

{\large Reaktionsgleichung:}
\vspace{0.5cm}

{\Huge \lsg{2 Ba + O$_2$ $\rightarrow$ 2 BaO}}
\vspace{0.5cm}

{\small Kontrolle: 2 Ba, 2 O links = 2 Ba, 2 O rechts \checkmark}
\end{center}
\end{frame}

\begin{frame}
\frametitle{1. Bariumoxid – Funfact}
\begin{center}
\vspace{0.3cm}
{\Huge \textcolor{chemie}{BaO}}
\vspace{0.8cm}

{\Large\textbf{Anwendung: Feuerwerk!}}
\vspace{0.5cm}

Bariumsalze erzeugen die typische\\[0.3cm]
{\Huge \textcolor{green!50!black}{\textbf{grüne Flammenfarbe}}}
\vspace{0.8cm}

\textcolor{akzent}{\textbf{Funfact:}} Bis zu \textbf{500 kg} Bariumverbindungen pro Show!
\end{center}
\end{frame}

% ═══════════════════════════════════════════════════════════════
% 2. NATRIUMOXID
% ═══════════════════════════════════════════════════════════════

\begin{frame}
\frametitle{2. Natriumoxid – Aufgabe}
\begin{block}{Gegeben}
Natrium (Na), Wertigkeit \textbf{I} \quad|\quad Sauerstoff (O), Wertigkeit \textbf{II}
\end{block}
\vspace{0.5cm}
\begin{block}{Aufgaben}
\begin{enumerate}
\item Oxidformel aufstellen: \luecke
\item Reaktionsgleichung (ausgeglichen): \quad \luecke{} + O$_2$ $\rightarrow$ \luecke
\end{enumerate}
\end{block}
\end{frame}

\begin{frame}
\frametitle{2. Natriumoxid – Lösung Schritt 1}
\begin{block}{Oxidformel}
Na (I) + O (II) $\rightarrow$ Kreuzregel: \textbf{kreuzen!} \hfill \lsg{Na$_2$O}
\end{block}
\vspace{0.3cm}
\begin{block}{Reaktionsgleichung (unausgeglichen)}
\lsg{Na + O$_2$ $\rightarrow$ Na$_2$O}\\[0.3cm]
Links: 1 Na, 2 O \quad Rechts: 2 Na, 1 O \quad $\rightarrow$ \textbf{nicht ausgeglichen!}
\end{block}
\end{frame}

\begin{frame}
\frametitle{2. Natriumoxid – Lösung komplett}
\begin{center}
\vspace{0.5cm}
{\large Oxidformel:} \quad {\Huge \lsg{Na$_2$O}}
\vspace{1cm}

{\large Reaktionsgleichung:}
\vspace{0.5cm}

{\Huge \lsg{4 Na + O$_2$ $\rightarrow$ 2 Na$_2$O}}
\vspace{0.5cm}

{\small Kontrolle: 4 Na, 2 O links = 4 Na, 2 O rechts \checkmark}
\end{center}
\end{frame}

\begin{frame}
\frametitle{2. Natriumoxid – Funfact}
\begin{center}
\vspace{0.3cm}
{\Huge \textcolor{chemie}{Na$_2$O}}
\vspace{0.8cm}

{\Large\textbf{Anwendung: Glasherstellung}}
\vspace{0.5cm}

Na$_2$O senkt den Schmelzpunkt von Glas
\vspace{0.8cm}

\textcolor{akzent}{\textbf{Funfact:}} Jede Glasflasche enthält ca. \textbf{15\%} Natriumoxid –\\
sonst müsste man bei über \textbf{2000°C} schmelzen!
\end{center}
\end{frame}

% ═══════════════════════════════════════════════════════════════
% 3. CHROMOXID
% ═══════════════════════════════════════════════════════════════

\begin{frame}
\frametitle{3. Chromoxid – Aufgabe}
\begin{block}{Gegeben}
Chrom (Cr), Wertigkeit \textbf{III} \quad|\quad Sauerstoff (O), Wertigkeit \textbf{II}
\end{block}
\vspace{0.5cm}
\begin{block}{Aufgaben}
\begin{enumerate}
\item Oxidformel aufstellen: \luecke
\item Reaktionsgleichung (ausgeglichen): \quad \luecke{} + O$_2$ $\rightarrow$ \luecke
\end{enumerate}
\end{block}
\end{frame}

\begin{frame}
\frametitle{3. Chromoxid – Lösung Schritt 1}
\begin{block}{Oxidformel}
Cr (III) + O (II) $\rightarrow$ Kreuzregel: \textbf{kreuzen!} \hfill \lsg{Cr$_2$O$_3$}
\end{block}
\vspace{0.3cm}
\begin{block}{Reaktionsgleichung (unausgeglichen)}
\lsg{Cr + O$_2$ $\rightarrow$ Cr$_2$O$_3$}\\[0.3cm]
Links: 1 Cr, 2 O \quad Rechts: 2 Cr, 3 O\\[0.2cm]
\textcolor{akzent}{Achtung:} 3 O ist ungerade, O$_2$ immer gerade!
\end{block}
\end{frame}

\begin{frame}
\frametitle{3. Chromoxid – Lösung komplett}
\begin{center}
\vspace{0.3cm}
{\large Oxidformel:} \quad {\Huge \lsg{Cr$_2$O$_3$}}
\vspace{0.8cm}

{\large Reaktionsgleichung:}
\vspace{0.4cm}

{\Huge \lsg{4 Cr + 3 O$_2$ $\rightarrow$ 2 Cr$_2$O$_3$}}
\vspace{0.4cm}

{\small Trick: Produkt verdoppeln $\rightarrow$ 6 O $\rightarrow$ 3 O$_2$}\\
{\small Kontrolle: 4 Cr, 6 O links = 4 Cr, 6 O rechts \checkmark}
\end{center}
\end{frame}

\begin{frame}
\frametitle{3. Chromoxid – Funfact}
\begin{center}
\vspace{0.3cm}
{\Huge \textcolor{chemie}{Cr$_2$O$_3$}}
\vspace{0.8cm}

{\Large\textbf{Pigment ``Chromgrün''}}
\vspace{0.5cm}

Farben, Keramik, Tarnfarbe beim Militär
\vspace{0.8cm}

\textcolor{akzent}{\textbf{Funfact:}} Der grüne Farbton auf \textbf{US-Dollarnoten}\\
enthält Chromoxid – deshalb ``\textbf{Greenbacks}''!
\end{center}
\end{frame}

% ═══════════════════════════════════════════════════════════════
% 4. TITANDIOXID
% ═══════════════════════════════════════════════════════════════

\begin{frame}
\frametitle{4. Titandioxid – Aufgabe}
\begin{block}{Gegeben}
Titan (Ti), Wertigkeit \textbf{IV} \quad|\quad Sauerstoff (O), Wertigkeit \textbf{II}
\end{block}
\vspace{0.5cm}
\begin{block}{Aufgaben}
\begin{enumerate}
\item Oxidformel aufstellen: \luecke
\item Reaktionsgleichung (ausgeglichen): \quad \luecke{} + O$_2$ $\rightarrow$ \luecke
\end{enumerate}
\end{block}
\end{frame}

\begin{frame}
\frametitle{4. Titandioxid – Lösung Schritt 1}
\begin{block}{Oxidformel}
Ti (IV) + O (II) $\rightarrow$ Kreuzregel: 4:2 $\rightarrow$ \textbf{kürzen: 1:2} \hfill \lsg{TiO$_2$}
\end{block}
\vspace{0.3cm}
\begin{block}{Reaktionsgleichung}
\lsg{Ti + O$_2$ $\rightarrow$ TiO$_2$}\\[0.3cm]
Links: 1 Ti, 2 O \quad Rechts: 1 Ti, 2 O \quad $\rightarrow$ \textbf{schon ausgeglichen!}
\end{block}
\end{frame}

\begin{frame}
\frametitle{4. Titandioxid – Lösung komplett}
\begin{center}
\vspace{0.5cm}
{\large Oxidformel:} \quad {\Huge \lsg{TiO$_2$}}
\vspace{1cm}

{\large Reaktionsgleichung:}
\vspace{0.5cm}

{\Huge \lsg{Ti + O$_2$ $\rightarrow$ TiO$_2$}}
\vspace{0.5cm}

{\small Einfach: Bereits ausgeglichen! \checkmark}
\end{center}
\end{frame}

\begin{frame}
\frametitle{4. Titandioxid – Funfact}
\begin{center}
\vspace{0.3cm}
{\Huge \textcolor{chemie}{TiO$_2$}}
\vspace{0.8cm}

{\Large\textbf{Weißpigment Nr. 1 weltweit!}}
\vspace{0.5cm}

Wandfarbe, Kunststoff, Sonnencreme, Zahnpasta
\vspace{0.8cm}

\textcolor{akzent}{\textbf{Funfact:}} \textbf{8 Millionen Tonnen} pro Jahr –\\
mehr als jedes andere Pigment!
\end{center}
\end{frame}

% ═══════════════════════════════════════════════════════════════
% 5. MANGANOXID
% ═══════════════════════════════════════════════════════════════

\begin{frame}
\frametitle{5. Manganoxid – Aufgabe}
\begin{block}{Gegeben}
Mangan (Mn), Wertigkeit \textbf{IV} \quad|\quad Sauerstoff (O), Wertigkeit \textbf{II}
\end{block}
\vspace{0.5cm}
\begin{block}{Aufgaben}
\begin{enumerate}
\item Oxidformel aufstellen: \luecke
\item Reaktionsgleichung (ausgeglichen): \quad \luecke{} + O$_2$ $\rightarrow$ \luecke
\end{enumerate}
\end{block}
\end{frame}

\begin{frame}
\frametitle{5. Manganoxid – Lösung Schritt 1}
\begin{block}{Oxidformel}
Mn (IV) + O (II) $\rightarrow$ Kreuzregel: 4:2 $\rightarrow$ \textbf{kürzen: 1:2} \hfill \lsg{MnO$_2$}
\end{block}
\vspace{0.3cm}
\begin{block}{Reaktionsgleichung}
\lsg{Mn + O$_2$ $\rightarrow$ MnO$_2$}\\[0.3cm]
Links: 1 Mn, 2 O \quad Rechts: 1 Mn, 2 O \quad $\rightarrow$ \textbf{schon ausgeglichen!}
\end{block}
\end{frame}

\begin{frame}
\frametitle{5. Manganoxid – Lösung komplett}
\begin{center}
\vspace{0.5cm}
{\large Oxidformel:} \quad {\Huge \lsg{MnO$_2$}}
\vspace{1cm}

{\large Reaktionsgleichung:}
\vspace{0.5cm}

{\Huge \lsg{Mn + O$_2$ $\rightarrow$ MnO$_2$}}
\vspace{0.5cm}

{\small Einfach: Bereits ausgeglichen! \checkmark}
\end{center}
\end{frame}

\begin{frame}
\frametitle{5. Manganoxid – Funfact}
\begin{center}
\vspace{0.3cm}
{\Huge \textcolor{chemie}{MnO$_2$}}
\vspace{0.8cm}

{\Large\textbf{``Braunstein'' in Batterien!}}
\vspace{0.5cm}

Zink-Kohle und Alkaline-Batterien
\vspace{0.8cm}

\textcolor{akzent}{\textbf{Funfact:}} Jede AA-Batterie enthält \textbf{8 g} Braunstein –\\
weltweit über \textbf{100.000 Tonnen} pro Jahr!
\end{center}
\end{frame}

% ═══════════════════════════════════════════════════════════════
% 6. ZINNOXID
% ═══════════════════════════════════════════════════════════════

\begin{frame}
\frametitle{6. Zinnoxid – Aufgabe}
\begin{block}{Gegeben}
Zinn (Sn), Wertigkeit \textbf{IV} \quad|\quad Sauerstoff (O), Wertigkeit \textbf{II}
\end{block}
\vspace{0.5cm}
\begin{block}{Aufgaben}
\begin{enumerate}
\item Oxidformel aufstellen: \luecke
\item Reaktionsgleichung (ausgeglichen): \quad \luecke{} + O$_2$ $\rightarrow$ \luecke
\end{enumerate}
\end{block}
\end{frame}

\begin{frame}
\frametitle{6. Zinnoxid – Lösung Schritt 1}
\begin{block}{Oxidformel}
Sn (IV) + O (II) $\rightarrow$ Kreuzregel: 4:2 $\rightarrow$ \textbf{kürzen: 1:2} \hfill \lsg{SnO$_2$}
\end{block}
\vspace{0.3cm}
\begin{block}{Reaktionsgleichung}
\lsg{Sn + O$_2$ $\rightarrow$ SnO$_2$}\\[0.3cm]
Links: 1 Sn, 2 O \quad Rechts: 1 Sn, 2 O \quad $\rightarrow$ \textbf{schon ausgeglichen!}
\end{block}
\end{frame}

\begin{frame}
\frametitle{6. Zinnoxid – Lösung komplett}
\begin{center}
\vspace{0.5cm}
{\large Oxidformel:} \quad {\Huge \lsg{SnO$_2$}}
\vspace{1cm}

{\large Reaktionsgleichung:}
\vspace{0.5cm}

{\Huge \lsg{Sn + O$_2$ $\rightarrow$ SnO$_2$}}
\vspace{0.5cm}

{\small Einfach: Bereits ausgeglichen! \checkmark}
\end{center}
\end{frame}

\begin{frame}
\frametitle{6. Zinnoxid – Funfact}
\begin{center}
\vspace{0.3cm}
{\Huge \textcolor{chemie}{SnO$_2$}}
\vspace{0.8cm}

{\Large\textbf{Keramik-Glasuren}}
\vspace{0.5cm}

Macht Glasuren weiß und deckend
\vspace{0.8cm}

\textcolor{akzent}{\textbf{Funfact:}} Die \textbf{Bronzezeit} begann, weil Menschen\\
Zinn aus SnO$_2$ gewannen – vor \textbf{5000 Jahren}!
\end{center}
\end{frame}

% ═══════════════════════════════════════════════════════════════
% 7. SILBEROXID
% ═══════════════════════════════════════════════════════════════

\begin{frame}
\frametitle{7. Silberoxid – Aufgabe}
\begin{block}{Gegeben}
Silber (Ag), Wertigkeit \textbf{I} \quad|\quad Sauerstoff (O), Wertigkeit \textbf{II}
\end{block}
\vspace{0.5cm}
\begin{block}{Aufgaben}
\begin{enumerate}
\item Oxidformel aufstellen: \luecke
\item Reaktionsgleichung (ausgeglichen): \quad \luecke{} + O$_2$ $\rightarrow$ \luecke
\end{enumerate}
\end{block}
\end{frame}

\begin{frame}
\frametitle{7. Silberoxid – Lösung Schritt 1}
\begin{block}{Oxidformel}
Ag (I) + O (II) $\rightarrow$ Kreuzregel: \textbf{kreuzen!} \hfill \lsg{Ag$_2$O}
\end{block}
\vspace{0.3cm}
\begin{block}{Reaktionsgleichung (unausgeglichen)}
\lsg{Ag + O$_2$ $\rightarrow$ Ag$_2$O}\\[0.3cm]
Links: 1 Ag, 2 O \quad Rechts: 2 Ag, 1 O \quad $\rightarrow$ \textbf{nicht ausgeglichen!}
\end{block}
\end{frame}

\begin{frame}
\frametitle{7. Silberoxid – Lösung komplett}
\begin{center}
\vspace{0.5cm}
{\large Oxidformel:} \quad {\Huge \lsg{Ag$_2$O}}
\vspace{1cm}

{\large Reaktionsgleichung:}
\vspace{0.5cm}

{\Huge \lsg{4 Ag + O$_2$ $\rightarrow$ 2 Ag$_2$O}}
\vspace{0.5cm}

{\small Kontrolle: 4 Ag, 2 O links = 4 Ag, 2 O rechts \checkmark}
\end{center}
\end{frame}

\begin{frame}
\frametitle{7. Silberoxid – Funfact}
\begin{center}
\vspace{0.3cm}
{\Huge \textcolor{chemie}{Ag$_2$O}}
\vspace{0.6cm}

{\Large\textbf{Das ``Anlaufen'' von Silber!}}
\vspace{0.4cm}

Schwarze Schicht auf altem Schmuck
\vspace{0.6cm}

\textcolor{akzent}{\textbf{Funfact:}} Die Reaktion lässt sich \textbf{umkehren}!\\
Erhitzt man Ag$_2$O $\rightarrow$ glänzendes Silber
\vspace{0.4cm}

{\Large $\rightarrow$ Das ist \textbf{Reduktion!}}
\end{center}
\end{frame}

% ═══════════════════════════════════════════════════════════════
% ÜBERLEITUNG
% ═══════════════════════════════════════════════════════════════

\begin{frame}
\frametitle{Oxidation $\leftrightarrow$ Reduktion}
\begin{center}
\vspace{0.3cm}
\begin{tikzpicture}[scale=0.95]
\node[draw=chemie, fill=hellgrau, rounded corners, minimum width=5.5cm, minimum height=2cm, align=center] at (-4,0) {
\textbf{Oxidation}\\[0.2cm]
4 Ag + O$_2$ $\rightarrow$ 2 Ag$_2$O\\[0.1cm]
{\small Sauerstoff-Aufnahme}
};
\draw[->, line width=2pt, akzent] (-0.7,0) -- (0.7,0);
\node[above] at (0,0.2) {\small Erhitzen};
\node[draw=akzent, fill=hellgrau, rounded corners, minimum width=5.5cm, minimum height=2cm, align=center] at (4,0) {
\textbf{\textcolor{akzent}{Reduktion}}\\[0.2cm]
2 Ag$_2$O $\rightarrow$ 4 Ag + O$_2$\\[0.1cm]
{\small Sauerstoff-Entzug}
};
\end{tikzpicture}
\vspace{0.8cm}

{\Large $\rightarrow$ \textbf{Nächstes Thema: Reduktion}}
\end{center}
\end{frame}

\end{document}
